\subsection{用户审核模块}

\paragraph*{业务流程}

\paragraph*{用例图}

\subsubsection{功能点——审核任务展示 （FR-YHSH-0001）}

\paragraph*{简述}

按照发布时间先后展示用户要求人工审核的图片条目，点击查看任务详情。支持按照学科、紧急程度等分类查看审核任务；支持详情页面展示已有的审核结果。

\paragraph*{前提条件}

由所有用户使用，用户必须已登录系统。

\paragraph*{主要流程}

业务流程为：用户登录后，进入“人工审核”页面，在任务列表中选择筛选条件（如学科：计算机视觉等），在界面上分类显示审核任务条目，点击任务条目查看详情。

\paragraph*{后继结果}

用户可对审核任务进行审核、评论等操作。

\subsubsection{功能点——用户审核（FR-YHSH-0002）}

\paragraph*{简述}

用户可对他人提交的人工审核申请的图片进行审查，提交鉴定结论及其理由。系统提供AI检测报告作为用户审核的参考。

\paragraph*{前提条件}

用户已进入审核任务详情页面。

输入信息：鉴定结论及其理由。

\paragraph*{主要流程}

业务流程为：用户进入审核任务详情页面后，点击“查看报告”按钮查看 AI 报告作为参考，选择鉴定结论（如“确认造假”、“未发现异常”等），填写文字理由（如“图1存在复制痕迹”），点击“提交审核”。

\paragraph*{后继结果}

输出信息： 发起审核用户收到结果通知。

\subsubsection{功能点——评论点赞（FR-YHSH-0003）}

\paragraph*{简述}

用户可对审核任务发表评论或点赞，促进学术交流。系统支持过滤敏感词，展示合规评论；支持按点赞数、发表时间排序评论。

\paragraph*{前提条件}

用户已进入审核任务详情页面。

输入信息：用户评论或点赞。

\paragraph*{主要流程}

业务流程为：用户进入审核任务详情页面后，可对该任务进行点赞操作；或点击评论输入框，发表评论（如“义眼顶针”）。

\paragraph*{后继结果}

输出信息：发起审核用户收到互动通知。

\subsubsection{功能点——用户举报（FR-YHSH-0004）}

\paragraph*{简述}

针对违规审核任务和评论，用户可选择对其进行举报，并说明情况交由管理端进行处理。

\paragraph*{前提条件}

用户发现违规内容（如审核图片违规、评论人身攻击等）。

输入信息：用户举报类型、理由。

\paragraph*{主要流程}

业务流程为：用户进入审核任务详情页面后，针对违规任务或者评论，点击“举报”按钮，选择举报类型（如“血腥暴力”等），填写理由说明（如“鬼图不是学术图片，令人生理不适”），提交举报后自动提交至管理端举报处理队列。

\paragraph*{后继结果}

输出信息：举报用户收到举报结果反馈。被举报者可能被禁言或限制使用等。

\subsubsection{功能点——论文人工审核（FR-YHSH-0005）}

\paragraph*{简述}

具备审核资格的用户对他人提交的**全篇论文**人工审核任务进行审查：在查看任务关联的原文摘要或脱敏片段、系统提供的 AI 与事实性子结论等参考信息基础上，提交个人鉴定结论与理由，供平台后续汇总。

\paragraph*{前提条件}

审核人已登录且具备参与论文人工审核的权限；任务处于“可审核”状态；审核人未与任务存在利益冲突等业务禁用情形（具体规则由平台配置）。

\paragraph*{主要流程}

\begin{enumerate}
    \item 审核人在任务列表或详情中打开“论文人工审核”任务。
    \item 系统展示论文相关信息及自动检测参考结果（段落高亮、概率、事实性子结论等，以后端下发为准）。
    \item 审核人选择结论类型（如“同意系统倾向”“存在重大异议”等，选项由产品定义），并填写文字理由与可选分项勾选。
    \item 审核人提交；系统校验必填项后写入审核记录，并通知任务发起人（若启用通知）。
\end{enumerate}

\paragraph*{后继结果}

审核记录持久化并进入汇总流程（见 FR-YHSH-0007）；重复提交策略（覆盖/禁止）以后端与产品约定为准。提交失败时提示原因。

\subsubsection{功能点——Review人工审核（FR-YHSH-0006）}

\paragraph*{简述}

具备审核资格的用户对**同行评审 Review** 类人工审核任务进行审查，结合系统对 Review 的自动检测参考（如模板化倾向等，以后端为准），提交鉴定结论与理由。

\paragraph*{前提条件}

审核人已登录且具备参与 Review 人工审核的权限；任务处于可审核状态。

\paragraph*{主要流程}

\begin{enumerate}
    \item 审核人打开 Review 审核任务详情，查看 Review 文本及自动检测参考信息。
    \item 审核人依据平台规范填写结论与理由，必要时引用系统标注的可疑片段。
    \item 审核人提交审核意见，系统保存并更新任务可见状态。
\end{enumerate}

\paragraph*{后继结果}

审核意见进入汇总；若任务需多名审核人，列表展示各轮次意见与状态。

\subsubsection{功能点——人工审核结果汇总（FR-YHSH-0007）}

\paragraph*{简述}

平台对同一人工审核任务下多名审核人或多轮次提交的结论进行**统计与汇总**，形成面向任务发起人的**最终人工审查结果**或阶段性汇总视图（如多数意见、分歧说明、待仲裁标记等），用于辅助判定内容真实性。汇总规则由后端实现，前端展示汇总结果与明细入口。

\paragraph*{前提条件}

任务下已存在至少一条审核记录，或已达到触发汇总规则的条件（如人数、截止时间，以后端为准）。

\paragraph*{主要流程}

\begin{enumerate}
    \item 任务发起人或被授权用户在任务详情中查看“人工审核汇总”区域。
    \item 前端请求汇总结果接口，展示最终结论摘要、各审核人结论对比表或时间线（字段由接口定义）。
    \item 用户可展开查看单人意见全文；若存在“未达成一致”等状态，展示平台提示与后续操作指引（如申请仲裁、补充材料，以实现为准）。
\end{enumerate}

\paragraph*{后继结果}

发起人获得清晰可读的汇总结论；尚无足够审核数据时展示“审核进行中”及当前已提交数量。

\paragraph*{补充说明}

汇总算法与终局判定权归属以后台业务规则为准；前端负责忠实呈现接口数据并处理空状态与权限不足场景。